\section{Giới thiệu đề tài}
\subsection{Lý do chọn đề tài}
Sự phát triển mạnh mẽ của thương mại điện tử trong những năm gần đây đã làm thay đổi toàn diện cách thức mua sắm và tiêu dùng của người dùng. Đặc biệt trong lĩnh vực bán lẻ điện thoại và thiết bị công nghệ, các doanh nghiệp lớn như CellphoneS, FPT Shop, Thế Giới Di Động đã xây dựng những nền tảng thương mại điện tử hiện đại, đáp ứng đa dạng nhu cầu từ việc tra cứu sản phẩm, so sánh giá, đến đặt hàng trực tuyến. Các hệ thống này được đầu tư bài bản cả về giao diện, trải nghiệm người dùng (UI/UX) lẫn hạ tầng kỹ thuật, mang lại sự tiện lợi và tin cậy cho khách hàng.

Tuy nhiên, khi khảo sát và trải nghiệm thực tế, nhóm nhận thấy rằng các nền tảng thương mại điện tử hiện nay vẫn chủ yếu hoạt động theo cơ chế tương tác truyền thống, nơi người dùng cần thao tác thủ công để tìm kiếm sản phẩm hoặc truy cập thông tin. Một số trang web đã tích hợp chatbot hoặc trợ lý ảo cơ bản, tuy nhiên các hệ thống này chỉ dừng lại ở mức trả lời câu hỏi hoặc hỗ trợ tư vấn đơn giản, chưa có khả năng thực hiện các tác vụ có mức độ tương tác sâu như đặt hàng, hủy đơn hàng, kiểm tra giỏ hàng, hoặc cung cấp thông tin mà khách hàng quan tâm một cách nhanh chóng và an toàn.

Trong bối cảnh trí tuệ nhân tạo (AI) và xử lý ngôn ngữ tự nhiên (Natural Language Processing – NLP) ngày càng phát triển, việc tích hợp Trợ lý Ảo thông minh vào hệ thống thương mại điện tử mở ra hướng tiếp cận mới – nơi người dùng có thể tương tác trực tiếp với hệ thống bằng ngôn ngữ tự nhiên, thay vì qua các bước điều hướng phức tạp. Điều này không chỉ giúp tăng tính cá nhân hóa trong trải nghiệm mua sắm, mà còn giảm tải cho nhân viên chăm sóc khách hàng, đồng thời nâng cao hiệu quả kinh doanh cho doanh nghiệp.

Xuất phát từ thực tế đó, nhóm quyết định lựa chọn và thực hiện đề tài “Hệ thống Thương mại Điện tử cho Bán lẻ Điện thoại tích hợp Trợ lý Ảo AI Agent” với mục tiêu xây dựng một mô hình thương mại điện tử thông minh, có khả năng tương tác tự nhiên, hiểu ngữ cảnh, hỗ trợ và thực hiện hành động thay cho người dùng. Đề tài không chỉ mang ý nghĩa thực tiễn trong việc ứng dụng công nghệ AI vào thương mại điện tử, mà còn góp phần nghiên cứu và phát triển các mô hình tương tác người – máy thế hệ mới.

\subsection{Mục tiêu đề tài}
Mục tiêu của đề tài là phát triển một nền tảng thương mại điện tử thông minh, ứng dụng công nghệ trí tuệ nhân tạo nhằm nâng cao trải nghiệm người dùng trong quá trình mua sắm trực tuyến. Cốt lõi của đề tài là phát triển một hệ thống có khả năng kết hợp giữa giao diện web thương mại điện tử truyền thống và Trợ lý Ảo sử dụng kỹ thuật xử lý ngôn ngữ tự nhiên (Natural Language Processing – NLP), cho phép người dùng tương tác trực tiếp với hệ thống bằng ngôn ngữ tự nhiên để thực hiện các thao tác như tìm kiếm, đặt hàng, kiểm tra giỏ hàng hay tra cứu thông tin.

Một trong những mục tiêu quan trọng của đề tài là tối ưu hóa trải nghiệm khách hàng bằng cách giảm thiểu thao tác thủ công và rút ngắn thời gian tìm kiếm sản phẩm. Thông qua Trợ lý Ảo AI Agent, người dùng có thể dễ dàng tiếp cận các sản phẩm phù hợp với nhu cầu cá nhân, đồng thời nhận được tư vấn nhanh chóng, chính xác và thân thiện hơn so với các phương thức tìm kiếm truyền thống. Hệ thống không chỉ dừng lại ở việc gợi ý thông tin, mà còn hướng tới khả năng thực thi hành động, giúp tự động hóa quy trình mua sắm trực tuyến.

Bên cạnh việc mang lại lợi ích cho người dùng, đề tài còn hướng đến việc nâng cao hiệu quả vận hành của doanh nghiệp bán lẻ thông qua tự động hóa các hoạt động chăm sóc khách hàng và xử lý đơn hàng. Việc tích hợp Trợ lý Ảo giúp giảm tải khối lượng công việc cho nhân viên hỗ trợ, đồng thời tạo ra kênh giao tiếp liên tục và linh hoạt giữa khách hàng và hệ thống.

Ở khía cạnh nghiên cứu, đề tài góp phần mở rộng ứng dụng của các mô hình AI trong thương mại điện tử, đặc biệt trong lĩnh vực tương tác người – máy (Human–Computer Interaction). Bằng cách áp dụng các kỹ thuật hiện đại về NLP và tích hợp hệ thống web, đề tài hướng tới việc chứng minh tiềm năng của Trợ lý Ảo trong việc nâng cao chất lượng dịch vụ và thúc đẩy chuyển đổi số trong ngành bán lẻ.

\subsection{Phạm vi đề tài}
Phạm vi của đề tài tập trung vào việc nghiên cứu, thiết kế và phát triển mô hình hệ thống thương mại điện tử thông minh, ứng dụng công nghệ trí tuệ nhân tạo nhằm nâng cao trải nghiệm người dùng trong quá trình mua sắm trực tuyến. Đề tài hướng đến việc kết hợp giữa nền tảng web thương mại điện tử truyền thống và Trợ lý Ảo AI Agent có khả năng hiểu, phản hồi và thực hiện hành động dựa trên ngôn ngữ tự nhiên.

Trước hết, đề tài bao gồm việc thiết kế và xây dựng kiến trúc tổng thể của hệ thống thương mại điện tử, với các thành phần chính như quản lý sản phẩm, giỏ hàng, đặt hàng, tra cứu thông tin đơn hàng và quản lý tài khoản người dùng. Hệ thống được định hướng hoạt động trên nền tảng web, đảm bảo khả năng mở rộng, an toàn và dễ sử dụng, đồng thời đáp ứng nhu cầu vận hành thực tế của doanh nghiệp trong lĩnh vực bán lẻ thiết bị di động.

Bên cạnh đó, đề tài tập trung vào phát triển mô hình Trợ lý Ảo AI Agent có khả năng xử lý ngôn ngữ tự nhiên (Natural Language Processing – NLP), giúp người dùng tương tác với hệ thống bằng ngôn ngữ tự nhiên thay vì thao tác truyền thống. Trợ lý Ảo không chỉ dừng lại ở việc tư vấn hoặc trả lời câu hỏi, mà còn được định hướng mở rộng khả năng thực hiện các hành động như đặt hàng, hủy đơn hàng hoặc kiểm tra giỏ hàng, tạo nên trải nghiệm tương tác liền mạch và thông minh hơn.

Đồng thời, đề tài cũng bao gồm thiết kế cơ sở dữ liệu và API nhằm phục vụ cho quá trình trao đổi thông tin giữa các thành phần hệ thống, đảm bảo tính toàn vẹn và hiệu năng cao trong quá trình truy xuất dữ liệu. Ngoài ra, giao diện người dùng (UI/UX) được chú trọng thiết kế theo hướng thân thiện, trực quan, phù hợp với thói quen sử dụng của khách hàng trong lĩnh vực thương mại điện tử.

Trong phạm vi giai đoạn đầu, nhóm tập trung vào các hoạt động phân tích yêu cầu, thiết kế hệ thống và xây dựng mô hình mô phỏng bao gồm: kiến trúc hệ thống, cơ sở dữ liệu (ERD, schema), thiết kế API, cùng với giao diện minh họa (mockup, wireframe). Việc triển khai hiện thực hệ thống, tích hợp AI Agent thực tế và kiểm thử tổng thể sẽ được thực hiện trong giai đoạn tiếp theo của đồ án.